\documentclass[10pt,a4paper]{amsart}
\usepackage[margin=19mm]{geometry}
\usepackage{amsmath,amssymb,mathtools}
\usepackage[scr=rsfs,cal=euler]{mathalfa}
\usepackage{xfrac}
\usepackage[unicode,hidelinks,bookmarks=false]{hyperref}
\newcommand{\ie}{i.e.\ }
\newcommand{\cf}{cf.\ }
\newcommand{\resp}{respectively\ }
\newcommand{\Subsection}{\subsection}
\providecommand{\nobreakdash}{\nobreak\hbox{-}}
\setlength{\emergencystretch}{2em}
\multlinegap0pt
\usepackage{bm}
\usepackage{bbm}
\usepackage{dsfont}
\usepackage{xspace}
\usepackage{cancel}
\usepackage{mathtools}
\usepackage[shortlabels]{enumitem}
\usepackage{amsthm}
\makeatletter
\@ifundefined{lemma}{\newtheorem{lemma}{Lemma}}{}
\@ifundefined{theorem}{\newtheorem{theorem}{Theorem}}{}
\makeatother
\DeclareMathOperator{\bigcomp}{\bigcirc}
\DeclareMathOperator{\prob}{\mathbb{P}}
\DeclareMathOperator{\rec}{flip}
\title[Singular functions obtained via random function iteration]{Singular functions obtained via random function iteration}
\author{Cristian Mitrea, Alef E. Sterk}
\date{Source: arXiv:2508.16327v1 (CC BY 4.0)}
\begin{document}
\maketitle

\noindent\emph{Source and licence.} Singular functions obtained via random function iteration. Version arXiv:2508.16327v1 (CC BY 4.0), distributed under the Creative Commons Attribution 4.0 International licence, as recorded in the arXiv licence metadata for this exact version and shown on the abstract page https://arxiv.org/abs/2508.16327v1. Full rights notice: licenses/arxiv-2508.16327-CC-BY-4.0.txt.\par\smallskip

\noindent\textit{Editorial scope.} Counted unit: the Theorem whose source label is \texttt{theorem:cantor\_properties} in the article (source0.tex lines 465--475, source label \texttt{theorem:cantor\_properties}), delivered together with the article's own complete original proof of it (source0.tex lines 477--555, one \texttt{proof} environment of 79 lines). No mathematics is written, repaired or re-derived: statement and proof are the article's own text line for line. Required context retained uncounted: lemma:sigma1 (lines 136--144); eq:def\_sequence (lines 249--252); lemma:properties2 (lines 315--322); lemma:properties3 (lines 378--387). Every definition, display and label of those lines is retained unchanged. Omitted from this package and not redistributed: 1--135, 145--248, 253--314, 323--377, 388--464, 476--476, 556--887. Nothing inside a retained excerpt is omitted or altered: every retained body is delivered exactly as the article's lines are written, so no omission receipt of any kind accompanies this package. Citations: the retained excerpts carry no citation command at all, so no bibliography entry of the article is transcribed and no reference is redistributed. Reference adaptation: none was needed and none was made; every reference inside the delivered unit resolves inside it, so no unresolved reference or citation is produced. Printed slips kept literal: no printed word, symbol, hypothesis or number of the counted statement or of its proof was altered. Layout and class: the article's own class is \texttt{article} with packages including amsmath, amssymb, amsthm, enumerate, graphicx, geometry, parskip, color, lineno, endfloat; the wrapper is the lane's pinned \texttt{amsart} editorial wrapper at 10pt on A4, which restates the theorem-like environments the retained text needs and adds the article's own macro definitions that the retained text needs (\texttt{\textbackslash bigcomp}, \texttt{\textbackslash prob}, \texttt{\textbackslash rec}), with extra wrapper packages (bm, bbm, dsfont, xspace, cancel, mathtools, enumitem). The delivered numbering is the unit's own. Assets: no retained span contains a figure environment, a graphics-inclusion command, a PDF-page inclusion, a table or a TikZ picture; no image file is copied into this package, the package declares no asset dependency and no image file is redistributed with it. Licence: version arXiv:2508.16327v1 of this article is distributed under the Creative Commons Attribution 4.0 International licence, as recorded in the arXiv licence metadata for this exact version and shown on the abstract page https://arxiv.org/abs/2508.16327v1. A full rights notice is delivered at licenses/arxiv-2508.16327-CC-BY-4.0.txt. Characters: every retained excerpt is pure ASCII, so the delivered engine, XeLaTeX with its default Latin Modern text font, reports no missing character.
\begin{lemma}
\label{lemma:sigma1}
For any measurable set $A \subseteq \Omega$ we have:
\begin{enumerate}[(i)]
\item $\prob_p(\sigma^{-1}(A) \cap [0]) = p\prob_p(A)$; 
\item $\prob_p(\sigma^{-1}(A) \cap [1]) = q\prob_p(A)$;
\item $\prob_p(\rec(A)) = \prob_q(A)$.
\end{enumerate}
\end{lemma}

\begin{equation}
\label{eq:def_sequence}
x_n = \bigcomp_{i=1}^n f_{\omega_i}(x_0) = (f_{\omega_n} \circ \dots \circ f_{\omega_2} \circ f_{\omega_1})(x_0), \quad n \in \mathbb{N}.
\end{equation}

\begin{lemma}
\label{lemma:properties2}
For any fixed $x_0, y_0 \in \mathbb{R}$ and $\omega \in \Omega$ the sequences $(x_n)$ and $(y_n)$ defined in equation~\eqref{eq:def_sequence} satisfy the following properties:
\begin{enumerate}[(i)]
\item If $x_0 \leq y_0$, then $x_n \leq y_n$ for all $n \in \mathbb{N}_0$;
\item Assume that (A4) holds. If $y_0 = 1-x_0$ but in the definition of $(y_n)$ the symbol sequence $\omega$ is replaced by the flipped sequence $\overline{\omega}$, then $y_n = 1-x_n$ for all $n \in \mathbb{N}$. In particular, $x_n \to \pm\infty$ if and only if $y_n \to \mp \infty$.
\end{enumerate}
\end{lemma}

\begin{lemma}
\label{lemma:properties3}
For any fixed $x_0 \in \mathbb{R}$ the sequence $(x_n)$ satisfies the following properties:
\begin{enumerate}[(i)]
\item If $x_0 \leq 0$, then $\prob_p(\{\omega \in \Omega \,:\, x_n \to \infty\}) = 0$;
\item If $x_0 \geq 1$, then $\prob_p(\{\omega \in \Omega \,:\, x_n \to \infty\}) = 1$;
\item $\prob_p(\{\omega \in \Omega \,:\, (x_n) \subseteq [0,1]\}) = 0$;
\item $\prob_p(\{\omega \in \Omega \,:\, (x_n) \to -\infty\}) = 1-\prob_p(\{\omega \in \Omega \,:\, (x_n) \to \infty\})$.
\end{enumerate}
\end{lemma}

\begin{theorem}
\label{theorem:cantor_properties}
The function $F_{p}$ has the following properties:
\begin{enumerate}[(i)]
\item $F_{p}(x) = 0$ for all $x \leq 0$ and $F_{p}(x) = 1$ for all $x \geq 1$;
\item $F_{p}(x) \leq F_p(y)$ whenever $x \leq y$;
\item $F_p(x) = pF_p(f_0(x)) + qF_p(f_1(x))$ for all $x \in \mathbb{R}$;
\item $F_p(x) = p$ for all $x \in [a,b]$;
\item $F_{p}(1-x) = 1-F_{q}(x)$ for all $x \in \mathbb{R}$ whenever assumption (A4) holds.
\end{enumerate}
\end{theorem}

\begin{proof}
(i) This follows directly from Lemma \ref{lemma:properties3}(i) and (ii).

(ii) If $x \leq y$, then Lemma \ref{lemma:properties2}(i) implies that
\[
\{ \omega \in \Omega \,:\, \bigcomp_{i=1}^n f_{\omega_i}(x) \to \infty \}
	\subseteq
	\{ \omega \in \Omega \,:\, \bigcomp_{i=1}^n f_{\omega_i}(y) \to \infty \},
\]
from which the claim follows.

(iii) Define the sets
\[
A = \{ \omega \in \Omega \,:\, \bigcomp_{i=1}^n f_{\omega_i}(x) \to \infty \}
\quad\text{and}\quad
B_k = \{ \omega \in \Omega \,:\, \bigcomp_{i=1}^n f_{\omega_i}(f_k(x)) \to \infty \},
\]
where $k \in \{0,1\}$. We claim that these sets are related as follows:
\[
A \cap [0] = \sigma^{-1}(B_0) \cap [0]
\quad\text{and}\quad
A \cap [1] = \sigma^{-1}(B_1) \cap [1].
\]
Indeed, we have the following equivalences:
\[
\begin{split}
(\omega_1,\omega_2,\omega_3,\dots) \in \sigma^{-1}(B_0) \cap [0]
	& \quad\Leftrightarrow\quad \omega_1 = 0 \text{ and } (\omega_2,\omega_3,\omega_4,\dots) \in B_0 \\
	& \quad\Leftrightarrow\quad \omega_1 = 0 \text{ and } \bigcomp_{i=2}^{n+1} f_{\omega_i}(f_0(x)) \to \infty \\
	& \quad\Leftrightarrow\quad \omega_1 = 0 \text{ and } \bigcomp_{i=1}^{n} f_{\omega_i}(x) \to \infty \\
	& \quad\Leftrightarrow\quad (\omega_1,\omega_2,\omega_3,\dots) \in A \cap [0].
\end{split}
\]
The other equality follows in a similar manner. Finally, using Lemma \ref{lemma:sigma1} gives
\[
\begin{split}
F_p(x)
	& = \prob_p(A) \\
	& = \prob_p(A \cap [0]) + \prob_p(A \cap [1]) \\
	& = \prob_p(\sigma^{-1}(B_0) \cap [0]) + \prob_p(\sigma^{-1}(B_1) \cap [1]) \\
	& = p\prob_p(B_0) + q\prob_p(B_1) \\
	& = pF_p(f_0(x)) + qF_p(f_1(x)),
\end{split}
\]
as desired.

(iv) Since
\[
f_0(a)=1, \quad
f_1(a)<0, \quad
f_0(b)>1, \quad
f_1(b)=0,
\]
it follows from statements (i) and (iii) that $F_p(a) = F_p(b) = p$. Statement (ii) implies that $F_p(x)=p$ for all $x \in [a,b]$.

(v) Under assumption (A4) we can apply Lemma \ref{lemma:properties2}(ii), which gives
\[
\begin{split}
F_{p}(1-x)
	& = \prob_p(\{ \omega \in \Omega \,:\, \bigcomp_{i=1}^n f_{\omega_i}(1-x) \to \infty \}) \\
	& = \prob_p(\{ \omega \in \Omega \,:\, \bigcomp_{i=1}^n f_{\overline{\omega}_i}(x) \to -\infty \}).
\end{split}
\]
Using Lemma \ref{lemma:properties3}(iv) gives
\[
F_{p}(1-x)
	= 1-\prob_p(\{ \omega \in \Omega \,:\, \bigcomp_{i=1}^n f_{\overline{\omega}_i}(x) \to \infty \}).
\]
Finally, using Lemma \ref{lemma:sigma1} gives
\[
\begin{split}
F_{p}(1-x)
	& = 1-\prob_p(\rec(\{ \omega \in \Omega \,:\, \bigcomp_{i=1}^n f_{\omega_i}(x) \to \infty \})) \\
	& = 1-\prob_q(\{ \omega \in \Omega \,:\, \bigcomp_{i=1}^n f_{\omega_i}(x) \to \infty \}) \\
	& = 1-F_{q}(x),
\end{split}
\]
as desired.
\end{proof}

\end{document}
